\section{Actual old-model points obtained from cubic type-II points}
\label{sec:cubic-typeii-point}

All source paths in this section are relative to
\path{src/Development/AutomorphicArithmetic/AutomorphicVarieties/}
at the pinned commit in AC-K29.  The conclusions concern actual
specialized points.  They do not identify the cubic deformation ring
with the old ring, and do not settle the mixed tangent-image condition
of AC-K29-tangent.

Let $D=\operatorname{Dat}(\Omega,K,c,e,\psi)$ be literal old source
data, write $O=c.O$, $\kappa=\overline{\mathbb Q}_\ell$, and put
\[
 \Gamma=G_{\mathbb Q},\quad H=G_K=\ker\eta,\quad
 \gamma=\kappa_\ell^e\psi,\quad c_c^2=1,\quad
 \eta(c_c)=-1,\quad \gamma(c_c)=1.                         \tag{GP.1}
\]
Here the symbol $\kappa_\ell$ denotes the cyclotomic character,
whereas $\kappa$ denotes the coefficient field.  The last identity
is the parity condition in the old datum.  Retain every other
condition in its \texttt{GD}; in particular the use below of an old
datum does not remove its conductor, inertia, auxiliary-field, or
genericity restrictions.
Let $\chi:H\to\mu_3$ satisfy $\chi^{c_c}=\chi^{-1}$, and fix
$j\in\{1,2\}$.  Define
\[
 q(hc_c^a)=\chi(h),\qquad h\in H,\quad a\in\{0,1\}.
 \quad q(gg')=q(g)q(g')^{\eta(g)},\quad q(c_c)=1.           \tag{GP.2}
\]
Use the source semidirect group
$\mathrm{GG}_3=\mathrm{GL}_3\rtimes\{\pm1\}$, with
$\theta(A)=J A^{-t}J$, and write $a(g)$ for the matrix component.
The cubic residual seed in this section is explicitly
\[
 \overline r^{\chi,j}(g).\mathrm{mat}
  =\operatorname{diag}\bigl(\overline q(g)^j\overline\gamma(g),
           1,\overline q(g)^j\overline\gamma(g)^{-1}\bigr),
 \qquad \overline r^{\chi,j}(g).\mathrm{mult}=\eta(g).
                                                               \tag{GP.3}
\]
The anti-character identity in (GP.2) makes this a homomorphism
with the displayed semidirect law.

\begin{lemma}[AC-K29-kill: residual-kernel killing survives endpoint untwisting]
\label{AC-K29-kill}
Suppose $\ell\ne3$ and $r^{\chi,j}:\Gamma\to\mathrm{GG}_3(\kappa)$
is continuous, integral, and strictly congruent to (GP.3).
Assume $a(c_c)=I$, its restriction to $H$ has endpoint block
$\{0,2\}$ and middle character $1$, and it kills the literal
modified set
\[
 \mathcal K_{\chi,j}=
 \{g:\ g\in I_v\text{ for some }v\nmid\ell,
                   \ \overline r^{\chi,j}(g)=1\}.            \tag{GP.4}
\]
Then the formula
\[
 \widetilde a(g)
   =\operatorname{diag}(q(g)^{-j},1,q(g)^{-j})a(g),
 \qquad \widetilde r(g).\mathrm{mult}=\eta(g)                \tag{GP.5}
\]
defines a continuous integral homomorphism.  It is strictly
congruent to the old $r_0(\eta,\gamma)$, is normalized for every
old marked element $g_1$, and kills the old set
\[
 \mathcal K_0=\{g:\ g\in I_v\text{ for some }v\nmid\ell,
                         \ \overline r_0(g)=1\}.             \tag{GP.6}
\]
In particular this implication does not require $\psi$ to be
unramified at the conductor of $\chi$.
\end{lemma}
\begin{proof}
For $g=hc_c$ the identity $a(c_c)=I$ gives $a(g)=a(h)$.
Thus every $a(g)$ is block diagonal for the endpoint/middle
decomposition, with middle entry $1$.  Put
$T(g)=\operatorname{diag}(q(g)^{-j},1,q(g)^{-j})$.
Every $T(g')$ commutes with every $a(g)$, and
$\theta(T(g'))=T(g')^{-1}$.  The identity (GP.2) now gives
\[
 T(gg')a(gg')
   =T(g)a(g)\,\theta_{\eta(g)}(T(g')a(g')),
                                                               \tag{GP.7}
\]
which proves the homomorphism assertion.  The entries of $T$ are
locally constant units, so continuity and integrality are preserved.
Multiplication by $T$ cancels the $q^j$ factors in (GP.3).
The strict residual congruence therefore becomes the old one,
globally on $\Gamma$.

The literal \texttt{Normalised} definition in
\path{UnitaryThreeGroups/UnitaryThreeReflectionRight.lean}
requires $a(c_c)=I$ and the $(0,1)$ and $(1,0)$ entries of
\[
 z(c_c,g_1)=a(g_1)-a(c_cg_1^{-1}c_c)                        \tag{GP.8}
\]
to vanish.  They vanish for every $g_1$ because both matrices
are block diagonal.  Also $q(c_c)=1$, so normalization is
preserved without imposing a condition on $\chi(g_1)$.

Take $g\in\mathcal K_0$.  Its old residual multiplier is $1$,
so $g\in H$, and its old residual first diagonal entry gives
$\overline\gamma(g)=1$.  Hence
\[
 \overline r^{\chi,j}(g)
     =\operatorname{diag}(\overline\chi(g)^j I_2,1)
       \quad\hbox{in endpoint/middle order}.                  \tag{GP.9}
\]
The element $g^3$ belongs to the same inertia group and its
modified residual image is $1$.  By (GP.4),
$r^{\chi,j}(g)^3=1$.  If $C$ is its endpoint matrix, then
$C^3=I$ and $C\equiv\chi(g)^jI\pmod{\mathfrak m_\kappa}$.
Choose a finite extension containing its entries and $\mu_3$.
Because $\ell\ne3$, the spectral idempotents
\[
       e_\zeta={1\over3}\sum_{a=0}^2\zeta^{-a}C^a,
                       \qquad \zeta\in\mu_3,                 \tag{GP.10}
\]
are integral, mutually orthogonal, and sum to $I$.
Their reductions are $I$ for $\zeta=\chi(g)^j$ and $0$
for the other two roots.  An integral idempotent whose matrix
reduces to $0$ is $0$: its image is a finite direct summand
of the free valuation-ring module, and Nakayama applies.
Thus $C=\chi(g)^jI$.  Equation (GP.5) gives
$\widetilde r(g)=1$.  This proves every field claimed.
\end{proof}

\begin{proposition}[AC-K29-point: actual factorization through the old source ring]
\label{AC-K29-point}
In the preceding setting assume the old datum $D$ and its
source coefficient embedding are fixed.  Suppose the values of
$r^{\chi,j}$ lie in a finite extension of $\mathbb Q_\ell$.
For some old $U:D.\operatorname{ULoc}$, assume there is an actual
local map
\[
 f_U:U.R\longrightarrow_O\kappa,
 \quad\operatorname{IsLocPt}(f_U),\quad
 D.\mathrm{loc}.I_t(U)\subseteq\ker f_U,
 \quad\operatorname{push}(f_U,U.r)
                   =\widetilde r\circ\iota.                 \tag{GP.11}
\]
Then $\widetilde r$ satisfies every literal field of
$D.\operatorname{IsPt}(U,g_1,\widetilde r)$ for every old
marked $g_1$.

If $S:D.\operatorname{Ring}(U,g_1)$ is any actual source
universal ring, there is an $O$-algebra map
$f_S:S.R\to\kappa$ with
\[
 \operatorname{IsLocPt}(f_S),\qquad
 \operatorname{push}(f_S,S.r)=\widetilde r.                  \tag{GP.12}
\]
Its restriction to $\operatorname{entAlg}_O(S.r)$ supplies
the exact \texttt{IsPush} point required by the source
\texttt{LemWAlt}.

The finite-coefficient hypothesis is automatic for an actual
$\operatorname{IsLocPt}$ specialization of a complete Noetherian
local $O$-algebra satisfying the source \texttt{IsCO}, by
\path{PadicHodgeTheory/PtFinFinite.lean}, \texttt{Loc.ptFin}.
This observation applies only after that actual specialization
and its ring hypotheses have been supplied.
\end{proposition}
\begin{proof}
\proofstep{1. The source fields.}
The seven fields of \texttt{Dat.IsPt}, in
\path{UnitaryThreeGroups/RhoKApply.lean}, are the multiplier,
continuity, integrality, strict residual congruence,
normalization, killing of \texttt{D.Kill}, and the local map
(GP.11).  The preceding lemma proves the first six, and (GP.11)
is precisely the seventh.  This proves the first assertion;
it remains to justify the stronger, actual ring factorization.

\proofstep{2. A coefficient ring with the correct residue field.}
Apply the source \texttt{Loc.ptFin} to the constant dual-number
lift of $f_U$.  Its hypotheses hold for the old $U.R$ and
coefficient $O$; it puts the entire image of $f_U$ in one finite
extension.  Enlarge this extension to a finite $E/\mathbb Q_\ell$
containing $\operatorname{Frac}O$, the values of $q$, and all
matrix entries of $r^{\chi,j}$.  Write $O_E$ and $\mathfrak m_E$
for its valuation ring and maximal ideal, and set
\[
        A=O+\mathfrak m_E\subset O_E.                        \tag{GP.13}
\]
This is the inverse image of the original residue field $k_O$
inside $k_E$.  It is a finite $O$-module: it is an $O$-submodule
of the finite module $O_E$.  It is a local ring with maximal
ideal $\mathfrak m_E$ and residue field $k_O$.  Indeed, an
element with nonzero residue in $k_O$ has inverse in $O_E$
whose residue still lies in $k_O$, hence inverse in $A$.
It is Noetherian because it is finite over the Noetherian ring
$O$.

For a uniformizer $\pi$ of $O$, write $a=v_E(\pi)$ in
integer normalization.  The powers of the maximal ideal of
$A$ are the sets $\mathfrak m_E^n$: a product of $n$ elements
of $\mathfrak m_E$ gives one inclusion, and
$\pi_E^n u=(\pi_Eu)\pi_E^{n-1}$ gives the other.
Consequently
\[
       \mathfrak m_A^{a+1}\subseteq\pi A
                      \subseteq\mathfrak m_A.                \tag{GP.14}
\]
The $\pi$-adic and $\mathfrak m_A$-adic topologies agree.
Since $A$ is a finite module over the complete DVR $O$, it is
complete.  The map $O\to A$ is local and its residue map is
surjective.  These are exactly the fields of \texttt{IsCO O A}
in \path{CommutativeAlgebra/ResMapBijective.lean}.
The construction of $A$ is necessary here: $O_E$ itself may
have a larger residue field and need not satisfy that record.

Every entry of $\widetilde r$ lies in $A$, because its reduction
is an entry of the old seed over $O$.  The inverse matrices also
lie in $A$, either by applying the same assertion to group
inverses and $\theta$, or by the adjugate formula and the unit
residual determinant.  Thus $\widetilde r$ is a homomorphism
to $\mathrm{GG}_3(A)$.  Its continuity is the required adic
continuity, since the topologies in (GP.14) agree with the
valuation topology restricted to $A$.  Its strict residual
congruence is exactly the \texttt{GLift.red} identity.

For every $x\in U.R$, the field \texttt{IsCO.res\_surj}
gives $b\in O$ with $x-b\in\mathfrak m_{U.R}$.
By \texttt{IsLocPt}, $f_U(x)-b\in\mathfrak m_E$.
Therefore $f_U(U.R)\subseteq A$.  Its factorization through
$A$ still kills $I_t(U)$ and yields the restricted lift.

\proofstep{3. Apply the literal universal property.}
By \texttt{GRepresents} in
\path{UnitaryThreeGroups/GRepresents.lean}, the universal ring
$S$ represents lifts with \texttt{NormCond} and
\texttt{BoxCond}; the latter consists exactly of killing the
specified set and admitting the local factorization which
kills $I_t(U)$.  The $A$-valued lift just constructed satisfies
all these conditions.  Hence there is an $O$-algebra map
$S.R\to A$ whose specialization of $S.r$ is $\widetilde r$.
Composing with $A\hookrightarrow\kappa$ gives $f_S$.

For completeness, every $O$-algebra map between two rings
satisfying \texttt{IsCO} is local.  If $x$ belongs to the
source maximal ideal but its image is a unit, choose an
$O$-unit $b$ with the same nonzero residue as that image.
Then $x-b$ is a unit in the source and its image lies in the
target maximal ideal, a contradiction.  Thus $f_S$ is
$\operatorname{IsLocPt}$.  This argument is also the complete
locality step in \texttt{res\_map} of
\texttt{ResMapBijective.lean}.  Restriction to the entry
algebra gives the matrix-entry equality in the definition of
\texttt{IsPush}; the multiplier is already $\eta$.
This proves (GP.12).

Finally, for a given actual local specialization $f:R\to\kappa$
of an \texttt{IsCO} ring, the map $F=f+0\epsilon$ satisfies the
hypotheses of \texttt{Loc.ptFin}.  That theorem puts all of
$f(R)$ in a single finite extension.  Matrix entries obtained
from $f$ therefore satisfy the finite-coefficient condition
used above.  No assertion about arbitrary abstract
$\kappa$-valued representations is needed.
\end{proof}

\begin{corollary}[AC-K29-typeII: old modularity and the original counting level]
\label{AC-K29-typeII}
Retain the old source \texttt{CohInp} and
\texttt{AutomorphyUnramifiedCoefficientsAlt}, its restriction
$e_O=1$, every old \texttt{GD} condition, and (GP.11) for each
cubic type-II point under consideration.  Suppose its endpoint
block on $G_K$ is irreducible, so that endpoint untwisting
preserves the source \texttt{TypeII} predicate.
For $\ell\notin\{3,269\}$, (GP.11) is supplied by
AC-G48-cubic when the point has the stated ordered ordinary
inertia flag, is crystalline on the source $H_L$, and its
endpoint determinant becomes $1$ after untwisting as in
AC-G45.  Those are explicit properties of the modified point;
no modified global deformation construction is presumed here.
Then its untwist is an actual member of the old rank-two
admissible family at level $N|D_K|$ used by AC-K29.  In
particular, the old source supplies a rational two-dimensional
$W$, an attached newform of weight $2e+1$ after the twist
$\kappa_\ell^{-e}$, the source \texttt{WitAQ}, and every
\texttt{Level} condition in \texttt{Adm2A}; these conclusions
are not inferred merely from conjugate self-duality.

For fixed $j$, endpoint untwisting is injective on isomorphism
classes of these $G_K$ representations.  At an inert rational
prime $p$ away from all specified conductors, where
$\chi(\operatorname{Frob}_w)=1$, it preserves their tested
Frobenius traces.  Thus this type-II subfamily uses the
unchanged fixed-$K$ counting input and the unchanged level
in AC-K29.  A family containing both $j=1,2$ can be partitioned
into those two subfamilies.
\end{corollary}
\begin{proof}
In the ordinary situation just specified, use the finite
coefficient ring (GP.13) and unrestricted local universality
to obtain the map $f_U$ specializing to the untwisted local
representation.  The locality argument above makes it an
\texttt{IsLocPt}.  AC-G48-cubic proves that it kills the
original $I_t$, which is exactly (GP.11).

The retained \texttt{AutomorphyUnramifiedCoefficientsAlt},
defined in
\path{ModularityArgument/ThresholdArguments/DenseThresholdCases/AutomorphyUnramifiedCoefficients.lean},
provides an old marked $g_1$ and \texttt{LemWAlt} at that
same $g_1$.  The normalization in AC-K29-kill holds for this
choice.  By \texttt{exists\_ring} in
\path{UnitaryThreeGroups/RingEx.lean}, the retained
\texttt{CohInp} supplies $S:D.\operatorname{Ring}(U,g_1)$.
AC-K29-point gives the actual \texttt{IsPush} point of $S$
required by \texttt{LemWAlt}.  Its middle line is $1$.
Consequently \texttt{LemWAlt} proves
$D.\operatorname{IsModularDense}(\iota,W)$ for every $W$
satisfying the literal \texttt{IsExt} predicate.
The alias \texttt{IsModularDense} is the definition
\texttt{IsModularA}, which includes both the newform and
\texttt{WitAQ}, in
\path{ModularityArgument/AutomorphyLifting/AutomorphyDescent/AutNewformW.lean}.

Now apply the complete theorem
\texttt{adm2B\_rhoK\_of\_typeII\_of\_modular\_extensions}
in
\path{ModularityArgument/ThresholdArguments/DenseThresholdCases/AdmissibleOfTypeII.lean}.
Its premises are precisely the retained \texttt{CohInp},
old marked datum, actual \texttt{IsPt}, \texttt{TypeII}, and
modularity of every \texttt{IsExt} just obtained.
Its proof constructs $W$ with determinant $\eta_K$ and checks
the inertia level away from $\ell$, the ordinary flag at
$\ell$, the no-link condition, and the determinant weight.
It concludes the rank-two admissibility at $N|D_K|$ with
the displayed middle character $1$.  This is the required
old family; in particular there is no extra factor $269$
introduced by the endpoint untwist.

An isomorphism between two original $G_K$ representations
$\chi^j W_1|_{G_K}\oplus1$ and
$\chi^j W_2|_{G_K}\oplus1$ preserves their unique
irreducible two-dimensional summands.  Untwisting them by
$\chi^{-j}$ proves that the old three-dimensional
representations are isomorphic, and the converse is
immediate.  At the test primes in the statement the scalar
twist is $1$, so equality of the tested traces is exact,
including inverse Frobenius.  Partitioning by $j$ avoids
any claim that images for the two distinct twists are
disjoint.  These facts prove the counting assertion.
\end{proof}

\begin{remark}[The remaining hypotheses after the point transfer]
Equation (GP.11) is a selected-ideal assertion.  AC-G48 closes
it for the stated crystalline, ordered ordinary type-II
points.  Crystallinity or a Weil--Deligne inertial type alone
has not been substituted for those hypotheses.  Pointwise
block untwisting does not
transport arbitrary off-block first-order lifts: those are
exactly the mixed directions retained in AC-K29-tangent.
Type-I automorphy, mixed tangent-image vanishing, integral
Selmer extraction, and exceptional coefficient primes remain
their own premises.  The case $\ell=3$ is excluded from
AC-K29-kill; it is not absorbed by the prime-to-$\ell$
idempotent argument.
\end{remark}

\paragraph{Geometric dictionary.}
On the common open arithmetic scheme obtained by inverting
$\ell N\,269\,D_K$, the character $\chi$ is the rank-one
coefficient of the finite cyclic cover already constructed.
For a stable lattice, each finite reduction of the endpoint
representation has the finite-cover interpretation of the
already audited SGA~1, Expos\'e~V, \S\S4 and~7.
Endpoint untwisting tensors the two-dimensional direct summand
with that rank-one coefficient's inverse and leaves the trivial
middle summand fixed.  It is an operation on this decomposed
point, not a tensor operation on an arbitrary three-dimensional
irreducible representation.  The residual-kernel calculation
proves factorization through the original permitted inertia
quotient; the ring $O+\mathfrak m_E$ handles the coefficient
deformation functor with its original residue field.
These categorical identifications use the earlier framework
readings only.  They supply neither an additional automorphy
theorem nor the mixed global Selmer vanishing.
